\chapter{Related work}
\label{ch:related_work}

Much of the previous work in this field has focused on developing and refining
rendering methods. We examine fundamental rendering techniques, early
optimizations using caching, and the more recent adoption of neural networks,
and show how all of these have influenced the latest state-of-the-art methods.

\section{Volume rendering}

To simplify volumetric rendering, early methods rendered volumes using
geometric primitives~\cite[Section~7.3]{kaufman_7_2005}. Most commonly, these
primitives were triangles and could be rendered using traditional techniques.
They approximate a surface within the volume, obtained using a segmentation
function. For one such function, the resulting surface is called an
\emph{isosurface}, and an example of an algorithm that generates and renders
it is the \emph{marching cubes} algorithm.

Accurate surface approximations may require a lot of space, but even then only
the surface is used and information from data within the volume is lost. The
answer to this problem was so-called \emph{direct volume rendering}
methods~\cite[Section~4]{kaufman_43_2000}. They generate an image directly
from 3D data. This can be achieved using object-order (how each data point
affects pixels), image-order (how all data points contribute to each pixel),
or domain-based techniques (transform 3D data into another domain and create
a projection).

One of the first image-order methods was \emph{ray marching}, which solves the
rendering equation by stepping along a ray. Later, the same problem was solved
stochastically instead using Monte Carlo
\emph{path tracing}~\cite{fong_production_2017}, which samples positions along
the ray rather than marching in fixed increments.

\section{Optimizations with caching}

Advanced rendering methods are computationally demanding, which has led to
the development of various optimization techniques. One such technique is
caching, first introduced for surface rendering Monte Carlo ray tracing by
Ward \emph{et~al.}~\cite{ward_ray_1988} and extended by
K\v{r}iv{\'a}nek \emph{et~al.}~\cite{krivanek_radiance_2008}. In this approach,
a cache stores radiance values. At each step, the system checks whether a cache
record exists. If it does, gradient interpolation is performed; otherwise, the
radiance is computed and stored in the cache.

Later, Khlebnikov \emph{et~al.}~\cite{khlebnikov_parallel_2014} applied caching
to volumetric rendering as well. Because volumetric rendering is usually
performed on graphics processors, traditional caching was not suitable
for this setting, so they had to run three tasks in parallel on the \ac{gpu}.
Their rendering procedure checks whether cached information is available for
each scattering event. If it is, irradiance is extrapolated; otherwise, new
entries are created. They also update the cache continuously to remove any
discontinuities.

\section{Neural networks}

Neural networks were first used to improve importance sampling for Monte Carlo
integration~\cite{muller_neural_2019}. In 2021, M\"uller \emph{et~al.} replaced
the traditional caching mechanism in surface rendering with a neural network
\cite{muller_real-time_2021}. They used a single network to predict radiance
values from spatio-directional coordinates, training it in real time. They cached
\emph{scattered radiance} because it is computationally the most expensive
component. The following year, M\"uller \emph{et~al.} introduced Multiresolution
Hash Encoding~\cite{muller_instant_2022}, which encodes position and direction
inputs into higher-dimensional representations to improve reconstruction quality
while still supporting various real-time tasks.

As with surface rendering, neural networks also began to appear in volumetric
rendering research. In Deep Scattering (2017), Kallweit
\emph{et~al.}~\cite{kallweit_deep_2017} used a model to render clouds.
The solution to the radiative transfer equation is a recursive relation that is
expensive to compute with Monte Carlo integration. They initially attempted to
predict the complete in-scattered radiance, but the model struggled with the
direct illumination radiance. Instead, they learned only the multi-scattered
paths, i.e., the \emph{indirect in-scattered radiance}, and used Monte Carlo
integration for the rest. One of the first uses of deep learning for volumetric
lighting is Deep Volumetric Ambient Occlusion (2020) by Engel and
Ropinski~\cite{engel_deep_2021}. They used a 3D encoder-decoder convolutional
neural network to predict ambient occlusion at a lower resolution. The volume
itself is rendered at full resolution, and the ambient occlusion is sampled
from the lower-resolution result.

\section{State of the art}

The latest state-of-the-art methods were introduced by Bauer \emph{et~al.}
In 2023, they presented Photon Field Networks~\cite{bauer_photon_2024},
replacing costly sampling steps with constant-time neural network inference to
achieve real-time volumetric global illumination. Their network predicted indirect
in-scatter radiance in the same way as Deep Scattering. They added the phase
function coefficient to position and direction inputs, for which they used
Multiresolution Hash Encoding. The model used fully connected layers and
predicted logarithm-encoded radiance for normalization. It was trained offline
using ground truth data obtained with photon mapping, which contains less
stochastic noise than path tracing. Then in 2025, they returned to caching
with GSCache~\cite{bauer_gscache_2026}. They used 3D Gaussian splatting functions
as a multi-level radiance cache. Each level is a point cloud of 3D Gaussians used
for paths of a specific length and is rasterized into images using Gaussian
splatting. In both cases, they improved performance over plain path tracing.
